\section{Conclusion}
\label{sec:conclusion}

We presented \method, a unified three-tier memory framework for
vision-language-action models. By separating memory along
\emph{time scale} (hot/cold/persistent) and \emph{modality}
(video/text/retrieval), we expose a coherent abstraction that
integrates with existing VLAs via small, documented patches rather
than wholesale rewrites.

Our initial release delivers two of the three tiers as a
zero-dependency Python package: a semi-structured language memory
schema that is both LLM-consumable and programmatically
inspectable, and a persistent cross-session store backed by
FAISS-ready retrieval. We demonstrate the framework's utility via
a Mem-0 integration that requires only two code edits, and we
propose MNEMO-Bench, a unified evaluation harness consolidating
four existing memory benchmarks. The package is open-sourced under
the MIT license, with 20 passing unit tests on CPU and a
reproducible quickstart and retrieval benchmark.

We see three directions of future work. First, implementing the
hot-memory tier with MEM-style video encoding and tactile fusion
will enable end-to-end policy training. Second, the joint training
objective proposed in Equation~\ref{eq:joint_loss} should be
validated empirically; preliminary infrastructure is in place.
Third, the MNEMO-Bench harness should be released with a public
leaderboard to accelerate community adoption of standardized memory
evaluation.

Memory has been the quiet bottleneck in robotic manipulation for
years; we hope \method\ lowers the cost of addressing it.
